\documentclass[12pt,letterpaper]{article}

% ── Packages ──────────────────────────────────────────────────────────────────
\usepackage[margin=1in]{geometry}
\usepackage{amsmath,amssymb}
\usepackage{graphicx}
\usepackage{hyperref}
\usepackage{xcolor}
\usepackage{booktabs}
\usepackage{array}
\usepackage{longtable}
\usepackage{enumitem}
\usepackage{fancyhdr}
\usepackage{titlesec}
\usepackage{setspace}
\usepackage{parskip}
\usepackage{microtype}
\usepackage{listings}
\usepackage[T1]{fontenc}
\usepackage{lmodern}

% ── Color Definitions ─────────────────────────────────────────────────────────
\definecolor{navyblue}{RGB}{31,56,100}
\definecolor{accentblue}{RGB}{46,116,181}
\definecolor{accentgreen}{RGB}{35,134,54}
\definecolor{accentorange}{RGB}{224,98,9}
\definecolor{lightgray}{RGB}{245,247,250}
\definecolor{codegray}{RGB}{100,110,120}
\definecolor{linkblue}{RGB}{56,139,253}

% ── Hyperref Setup ────────────────────────────────────────────────────────────
\hypersetup{
    colorlinks=true,
    linkcolor=accentblue,
    urlcolor=linkblue,
    citecolor=navyblue,
    pdfauthor={CS 599 -- DocuMind Project},
    pdftitle={DocuMind: A Multi-Module AI Platform for Regulatory Intelligence and Knowledge Synthesis},
}

% ── Section Formatting ────────────────────────────────────────────────────────
\titleformat{\section}
  {\normalfont\Large\bfseries\color{navyblue}}
  {\thesection.}{0.6em}{}[\color{navyblue}\titlerule]

\titleformat{\subsection}
  {\normalfont\large\bfseries\color{accentblue}}
  {\thesubsection}{0.6em}{}

\titleformat{\subsubsection}
  {\normalfont\normalsize\bfseries\color{accentorange}}
  {\thesubsubsection}{0.6em}{}

% ── Header / Footer ───────────────────────────────────────────────────────────
\pagestyle{fancy}
\fancyhf{}
\rhead{\small\color{codegray}CS 599 -- Applications of Large Language Models}
\lhead{\small\color{codegray}DocuMind Project Proposal}
\rfoot{\small\thepage}
\lfoot{\small\color{codegray}Northern Arizona University}
\renewcommand{\headrulewidth}{0.4pt}
\renewcommand{\footrulewidth}{0.4pt}

% ── Code Listing Style ────────────────────────────────────────────────────────
\lstset{
  backgroundcolor=\color{lightgray},
  basicstyle=\ttfamily\small,
  keywordstyle=\color{navyblue}\bfseries,
  commentstyle=\color{codegray}\itshape,
  stringstyle=\color{accentblue},
  breaklines=true,
  frame=single,
  rulecolor=\color{lightgray!60!black},
  framesep=4pt,
  xleftmargin=8pt,
  xrightmargin=8pt,
}

% ── Table Column Types ────────────────────────────────────────────────────────
\newcolumntype{L}[1]{>{\raggedright\arraybackslash}p{#1}}
\newcolumntype{C}[1]{>{\centering\arraybackslash}p{#1}}

% ─────────────────────────────────────────────────────────────────────────────
\begin{document}

% ── TITLE PAGE ────────────────────────────────────────────────────────────────
\begin{titlepage}
  \centering
  \vspace*{2.5cm}

  {\Huge\bfseries\color{navyblue} DocuMind\par}
  \vspace{0.6cm}
  {\Large\bfseries\color{accentblue}
    Production-Grade LLM Intelligence Platform \\
    for Regulatory Analysis and Knowledge Synthesis\par}
  \vspace{1.2cm}

  \noindent\rule{\textwidth}{1.5pt}
  \vspace{0.4cm}

  {\large\itshape
    Regulatory Change Analysis $\cdot$
    Multi-Agent Knowledge Synthesis $\cdot$
    Retrieval-Augmented Generation\par}
  \vspace{0.4cm}
  \noindent\rule{\textwidth}{0.4pt}

  \vspace{1.8cm}

  \begin{tabular}{rl}
    \textbf{Course:}      & CS 599 --- Contemporary Developments: Applications of LLMs \\[4pt]
    \textbf{Institution:} & Northern Arizona University \\[4pt]
    \textbf{Department:}  & Department of Computer Science \\[4pt]
    \textbf{Modules:}     & RegDelta $|$ SynapseIQ \\[4pt]
    \textbf{Repository:}  &
      \href{https://github.com/dp2426-NAU/LLM-Project}
           {\texttt{github.com/dp2426-NAU/LLM-Project}} \\
  \end{tabular}

  \vfill
  {\small\color{codegray} Version 1.0 $\cdot$ 2026}
\end{titlepage}

\tableofcontents
\newpage

% ─────────────────────────────────────────────────────────────────────────────
\section{Abstract}
% ─────────────────────────────────────────────────────────────────────────────

Organizations operating in regulated industries generate and consume enormous volumes of
structured and unstructured documents daily. Two critical knowledge-work bottlenecks remain
unsolved at enterprise scale: (1) compliance teams cannot efficiently propagate regulatory
amendments across internal policy inventories, and (2) knowledge workers lack automated
tools that synthesize multi-document corpora into rigorously evaluated artifacts.

This project presents \textbf{DocuMind}, a production-grade AI platform comprising two
complementary modules --- \textbf{RegDelta} and \textbf{SynapseIQ} --- both built on
Retrieval-Augmented Generation (RAG) and Large Language Model (LLM) inference.
RegDelta detects section-level regulatory changes, retrieves semantically matched internal
policies from a vector store, and generates risk-classified impact reports via an LLM
pipeline. SynapseIQ implements a four-agent sequential architecture
(Researcher $\to$ Critic $\to$ Synthesizer $\to$ Validator) that produces self-evaluated
knowledge syntheses with composite quality scoring.

Both modules are deployed as production FastAPI services with ChromaDB vector stores,
SQLite analytics databases, and Docker containerization. The platform demonstrates that
LLMs, when combined with retrieval grounding and multi-agent oversight, can significantly
reduce manual knowledge-work overhead in regulated domains.

% ─────────────────────────────────────────────────────────────────────────────
\section{Problem Statement}
% ─────────────────────────────────────────────────────────────────────────────

\subsection{Regulatory Change Propagation}

When regulatory bodies such as the SEC, FINRA, or Basel Committee amend rules, compliance
teams at financial and healthcare institutions must manually cross-reference hundreds of
internal policy documents to identify what requires updating. This process typically
requires \textbf{2--4 weeks} per major amendment, is highly error-prone under deadline
pressure, and produces no systematic audit trail linking identified policy gaps to specific
changed regulatory clauses. The problem scales linearly with the size of both the
regulation corpus and the organization's internal policy inventory.

\subsection{Research and Knowledge Synthesis at Scale}

Analysts, legal professionals, and researchers spend the majority of their cognitive effort
reading, extracting, and manually integrating information from heterogeneous document
corpora into coherent reports. Existing LLM-based tools --- document summarizers and Q\&A
chatbots --- retrieve relevant text but do not:
\begin{itemize}[leftmargin=1.5em]
  \item Reason about logical gaps and contradictions across multiple sources
  \item Apply structured critique before producing a final synthesis
  \item Self-evaluate output quality before delivery to a decision-maker
  \item Maintain a traceable reasoning chain from raw source documents to the final artifact
\end{itemize}

\subsection{Limitations of Existing Approaches}

\begin{table}[h!]
\centering
\begin{tabular}{L{4.5cm} L{9.5cm}}
\toprule
\textbf{Approach} & \textbf{Limitation} \\
\midrule
Simple RAG chatbots        & Returns raw chunks; no synthesized analysis; no quality gate \\
Document summarizers       & Single-document scope; no cross-document reasoning \\
Manual compliance review   & Non-scalable; no systematic gap detection; no audit trail \\
Generic LLM API calls      & No domain retrieval; hallucination risk without grounding \\
Rule-based diff tools      & No semantic understanding of regulatory significance \\
\bottomrule
\end{tabular}
\caption{Limitations of existing document intelligence approaches}
\end{table}

% ─────────────────────────────────────────────────────────────────────────────
\section{Objectives}
% ─────────────────────────────────────────────────────────────────────────────

\subsection{Module I --- RegDelta}

\begin{table}[h!]
\centering
\begin{tabular}{C{0.6cm} L{4.5cm} L{8.4cm}}
\toprule
\textbf{ID} & \textbf{Objective} & \textbf{Success Criterion} \\
\midrule
M1 & Section-level change detection   & ADDED/MODIFIED/REMOVED classification with similarity score \\
M2 & Semantic policy impact retrieval & Retrieve policies with cosine similarity $\geq 0.35$ \\
M3 & LLM-generated impact report      & JSON: risk\_level, executive\_summary, recommended\_actions \\
M4 & Version registry                 & SQLite tracks all regulation versions and reports with metadata \\
M5 & Production REST API              & FastAPI endpoints with OpenAPI docs, Pydantic validation, health check \\
\bottomrule
\end{tabular}
\caption{RegDelta module objectives and success criteria}
\end{table}

\subsection{Module II --- SynapseIQ}

\begin{table}[h!]
\centering
\begin{tabular}{C{0.6cm} L{4.5cm} L{8.4cm}}
\toprule
\textbf{ID} & \textbf{Objective} & \textbf{Success Criterion} \\
\midrule
S1 & 4-agent sequential pipeline  & Researcher$\to$Critic$\to$Synthesizer$\to$Validator with shared \texttt{AgentContext} \\
S2 & Self-evaluating output       & Composite score computed; PASS/CONDITIONAL\_PASS/FAIL verdict issued \\
S3 & Format-adaptive synthesis    & brief ($\sim$300w), standard ($\sim$800w), detailed ($\sim$2000w) modes \\
S4 & Agent observability          & Per-agent token count, latency, success recorded per session \\
S5 & Cost-bounded inference       & Per-agent token budgets enforced; USD cost tracked per query \\
\bottomrule
\end{tabular}
\caption{SynapseIQ module objectives and success criteria}
\end{table}

% ─────────────────────────────────────────────────────────────────────────────
\section{Methodology}
% ─────────────────────────────────────────────────────────────────────────────

\subsection{Retrieval-Augmented Generation (RAG)}

Both modules ground all LLM inference in retrieved document evidence rather than relying on
parametric knowledge. The RAG pipeline operates in two phases:

\subsubsection*{Offline Indexing}
Raw documents (PDF, TXT, Markdown) are parsed, cleaned, and split into overlapping
fixed-size chunks (512--800 tokens, 50--100 token overlap). Each chunk is embedded using
\texttt{sentence-transformers/all-MiniLM-L6-v2} (384-dimensional, L2-normalized dense
vectors) and upserted into ChromaDB with SHA-256-based document IDs for idempotent
re-ingestion.

\subsubsection*{Online Retrieval}
At query time, the user's query is embedded with the same model and passed to ChromaDB's
HNSW index for approximate nearest-neighbor (ANN) search. The top-$K$ chunks above a
minimum relevance threshold are returned as the LLM grounding context.

\subsection{Section-Level Differential Analysis (RegDelta)}

RegDelta implements a domain-specific diff engine at heading-level document granularity.
Documents are parsed using a compiled regex matching standard regulatory heading conventions
(\texttt{Section}, \texttt{Article}, \texttt{Rule}, \texttt{\S} prefixes). Python's
\texttt{difflib.SequenceMatcher} computes a similarity ratio in $[0, 1]$ for each matched
section pair. Sections are classified as \texttt{ADDED}, \texttt{REMOVED}, or
\texttt{MODIFIED}, and sorted by ascending similarity to surface the most significant
changes first.

\subsection{Multi-Agent Sequential Architecture (SynapseIQ)}

SynapseIQ implements a four-stage sequential agent pipeline, all sharing an
\texttt{AgentContext} dataclass that acts as a mutable state graph:

\begin{enumerate}[leftmargin=2em]
  \item \textbf{ResearcherAgent} --- Extracts structured findings, key concepts, and
        evidence gaps from retrieved chunks. YAML+Jinja2 prompt. Budget: 1,200 tokens.

  \item \textbf{CriticAgent} --- Identifies logical gaps, unsupported claims, and bias
        signals. Outputs structured JSON critique. Budget: 800 tokens.

  \item \textbf{SynthesizerAgent} --- Integrates findings and critique into a
        format-appropriate knowledge artifact (brief/standard/detailed).
        Budget: 1,800 tokens.

  \item \textbf{ValidatorAgent} --- Scores synthesis on four dimensions (coherence,
        relevance, factuality, completeness); issues a verdict. Budget: 600 tokens.
\end{enumerate}

All agents inherit from \texttt{BaseAgent}, which provides \texttt{\_timed\_llm\_call()}
(returns content, token count, and latency) and \texttt{\_safe\_run()} (with retry logic
and graceful degradation on failure).

\subsection{Evaluation Framework}

\subsubsection*{RegDelta}
Risk-level classification (\texttt{LOW}, \texttt{MEDIUM}, \texttt{HIGH}, \texttt{CRITICAL})
is determined by the LLM based on regulatory significance and affected policy count.

\subsubsection*{SynapseIQ}
A weighted composite evaluation score is computed before response delivery:

\begin{equation}
  \text{Composite} = 0.30 \times \text{coherence} + 0.35 \times \text{relevance}
                   + 0.25 \times \text{factuality} + 0.10 \times \text{completeness}
\end{equation}

Verdict thresholds:
\[
  \text{Verdict} =
  \begin{cases}
    \texttt{PASS}             & \text{if Composite} \geq 0.65 \\
    \texttt{CONDITIONAL\_PASS} & \text{if } 0.45 \leq \text{Composite} < 0.65 \\
    \texttt{FAIL}             & \text{if Composite} < 0.45
  \end{cases}
\]

% ─────────────────────────────────────────────────────────────────────────────
\section{System Architecture}
% ─────────────────────────────────────────────────────────────────────────────

\subsection{High-Level Layered Design}

\begin{lstlisting}[caption={DocuMind 5-layer system architecture}]
CLIENT LAYER     HTTP / REST  |  Swagger UI  |  Python SDK  |  CI/CD
       |
API GATEWAY      FastAPI / Uvicorn ASGI  |  Pydantic  |  OpenAPI 3.0
       |
DOMAIN MODULES   RegDelta  (Diff Engine -> ImpactAnalyzer -> Report)
                 SynapseIQ (Researcher -> Critic -> Synthesizer -> Validator)
       |
SHARED INFRA     ChromaDB  |  Embedding Engine  |  LLM Client  |  Config
       |
PERSISTENCE      SQLite (RegDelta versions + reports)
                 SQLite (SynapseIQ sessions + agent traces)
                 ChromaDB on-disk  |  OpenAI API / Ollama local
\end{lstlisting}

\subsection{RegDelta Data Flow}

\begin{lstlisting}[caption={RegDelta request lifecycle}]
POST /api/v1/analyze/delta
  -> VersionTracker.get_latest_two()         [fetch old + new text from SQLite]
  -> DiffEngine.compute_section_diffs()      [regex split + SequenceMatcher]
  -> ImpactAnalyzer._annotate_sections()     [LLM: 1 call per section, <=10]
  -> Retriever.retrieve_policies()           [ChromaDB ANN on policies collection]
  -> LLMGenerator.generate_json()            [full ImpactReport JSON]
  -> VersionTracker.save_report()            [persist to SQLite]
  <- ImpactReport.to_dict()                  [JSON response]
\end{lstlisting}

\subsection{SynapseIQ Pipeline Flow}

\begin{lstlisting}[caption={SynapseIQ multi-agent pipeline}]
POST /api/v1/synthesize
  -> CorpusStore.search(query, top_k=6)      [ChromaDB HNSW ANN search]
  -> ResearcherAgent._safe_run(context)      [YAML prompt -> LLM -> findings]
  -> CriticAgent._safe_run(context)          [YAML prompt -> LLM -> gaps]
  -> SynthesizerAgent._safe_run(context)     [YAML prompt -> LLM -> artifact]
  -> ValidatorAgent._safe_run(context)       [YAML prompt -> LLM -> scores]
  -> PipelineEvaluator.evaluate(context)     [composite score + verdict]
  -> PipelineTracer.end_session(...)         [SQLite analytics]
  <- SynthesisResponse JSON
\end{lstlisting}

% ─────────────────────────────────────────────────────────────────────────────
\section{Technologies Used}
% ─────────────────────────────────────────────────────────────────────────────

\begin{table}[h!]
\centering
\begin{tabular}{L{2.8cm} L{4.2cm} L{6.8cm}}
\toprule
\textbf{Category} & \textbf{Technology} & \textbf{Purpose} \\
\midrule
Language         & Python 3.11+           & Primary implementation language \\
Web Framework    & FastAPI + Uvicorn      & ASGI REST API server \\
LLM (cloud)      & OpenAI gpt-4o-mini     & Default inference backend \\
LLM (local)      & Ollama llama3          & Zero-cost local inference alternative \\
Vector Database  & ChromaDB 0.5           & Dense vector storage and ANN search \\
Embedding Model  & all-MiniLM-L6-v2      & 384-dim sentence embeddings \\
Persistence      & SQLite3                & Version tracking and session analytics \\
Data Validation  & Pydantic v2            & Request/response models and settings \\
Prompt System    & Jinja2 + YAML          & Versioned, hot-reloadable agent prompts \\
Diff Engine      & difflib (stdlib)       & Section-level text comparison \\
Containerization & Docker                 & Production deployment packaging \\
Testing          & pytest                 & Unit, integration, evaluation tests \\
\bottomrule
\end{tabular}
\caption{Technology stack for DocuMind}
\end{table}

% ─────────────────────────────────────────────────────────────────────────────
\section{Role of Large Language Models}
% ─────────────────────────────────────────────────────────────────────────────

LLMs serve as the \emph{reasoning core} of DocuMind. Critically, all LLM calls are supplied
with retrieved document context from ChromaDB, ensuring factual grounding and preventing
hallucination from parametric memory alone.

\subsection{LLM Usage in RegDelta}

\begin{table}[h!]
\centering
\begin{tabular}{L{3.5cm} C{1.5cm} L{3.5cm} L{4.5cm}}
\toprule
\textbf{Task} & \textbf{Calls} & \textbf{Input} & \textbf{Output} \\
\midrule
Section annotation    & $\leq$10 & Old + new text excerpt (400 chars each) & One-sentence significance note \\
Impact report gen.    & 1        & Changed sections + matched policy excerpts & JSON: risk\_level, summary, actions \\
\bottomrule
\end{tabular}
\caption{LLM calls in RegDelta per analysis request}
\end{table}

\subsection{LLM Usage in SynapseIQ}

\begin{table}[h!]
\centering
\begin{tabular}{L{3.5cm} L{4.5cm} C{2cm} L{3.2cm}}
\toprule
\textbf{Agent} & \textbf{Task} & \textbf{Budget} & \textbf{Output Format} \\
\midrule
ResearcherAgent  & Extract findings from retrieved chunks & 1,200 tok & core\_findings, key\_concepts, evidence\_gaps \\
CriticAgent      & Identify gaps, bias, unsupported claims & 800 tok   & Structured JSON critique \\
SynthesizerAgent & Produce knowledge artifact             & 1,800 tok & Formatted document \\
ValidatorAgent   & Score 4 dimensions; issue verdict      & 600 tok   & JSON scores + verdict \\
\bottomrule
\end{tabular}
\caption{LLM calls in SynapseIQ per synthesis request (4 agents, 4 calls)}
\end{table}

\subsection{Prompt Engineering Strategy}

SynapseIQ uses a YAML + Jinja2 prompt engineering system. Each agent's prompts are stored
in versioned YAML files containing \texttt{version}, \texttt{system}, \texttt{user}
(Jinja2 template), and \texttt{output\_schema} fields. The \texttt{PromptEngine} class
loads all templates at server startup and supports hot-reload via
\texttt{POST /api/v1/agents/prompts/reload}, enabling prompt iteration without restart.

\subsection{Cost Control}

Per-agent token budgets are enforced via \texttt{AGENT\_TOKEN\_BUDGETS} in
\texttt{constants.py}. A \texttt{CostReport} (estimated USD) is included in every
\texttt{/analytics} API response. Cost is estimated using:
\[
  \text{Cost} = \frac{T_\text{in} \times \$0.00015}{1000}
              + \frac{T_\text{out} \times \$0.00060}{1000}
  \quad \text{(gpt-4o-mini pricing)}
\]

% ─────────────────────────────────────────────────────────────────────────────
\section{Expected Outcomes}
% ─────────────────────────────────────────────────────────────────────────────

\begin{enumerate}[leftmargin=2em]
  \item \textbf{Automated regulatory compliance gap detection} --- RegDelta reduces manual
        policy review time from weeks to minutes by automatically identifying which internal
        policies are semantically affected by regulation amendments.

  \item \textbf{Rigorous multi-source knowledge synthesis} --- SynapseIQ produces
        quality-evaluated knowledge artifacts from multi-document corpora, with a traceable
        agent reasoning chain and a self-issued quality verdict before delivery.

  \item \textbf{Production-grade system design} --- Both modules demonstrate enterprise-
        quality software engineering: typed Python, dependency injection, structured logging,
        persistent storage, containerization, and comprehensive test coverage.

  \item \textbf{Transferable RAG + multi-agent patterns} --- The dual-corpus RAG,
        section-level diff, sequential agent state sharing, and YAML prompt management
        patterns are domain-agnostic and transferable to legal, medical, and scientific
        literature domains.

  \item \textbf{Observable LLM inference} --- Per-session agent traces, token counts,
        latency measurements, and cost estimates provide full observability into LLM usage,
        enabling optimization and debugging.
\end{enumerate}

% ─────────────────────────────────────────────────────────────────────────────
\section{Development Timeline}
% ─────────────────────────────────────────────────────────────────────────────

\begin{table}[h!]
\centering
\begin{tabular}{L{3.5cm} C{1.8cm} L{8.0cm}}
\toprule
\textbf{Phase} & \textbf{Duration} & \textbf{Deliverables} \\
\midrule
Phase 1 --- Infrastructure & Week 1 & ChromaDB, FastAPI skeleton, embedding pipeline, SQLite schema \\
Phase 2 --- RegDelta Core  & Week 2 & DiffEngine, ImpactAnalyzer, VectorStore, VersionTracker, REST API \\
Phase 3 --- SynapseIQ Core & Week 3 & All 4 agents, AgentContext, SynapseOrchestrator, PromptEngine \\
Phase 4 --- Evaluation     & Week 4 & PipelineEvaluator, metrics, cost tracker, analytics endpoints \\
Phase 5 --- Testing \& Docs & Week 5 & pytest suite (unit, integration, evaluation), README, proposal \\
Phase 6 --- Finalization   & Week 6 & Docker, GitHub Actions CI, performance tuning, final documentation \\
\bottomrule
\end{tabular}
\caption{Development timeline}
\end{table}

% ─────────────────────────────────────────────────────────────────────────────
\section{References}
% ─────────────────────────────────────────────────────────────────────────────

\begin{enumerate}[leftmargin=2em,label={[\arabic*]}]
  \item Lewis, P., et al. (2020).
        \textit{Retrieval-Augmented Generation for Knowledge-Intensive NLP Tasks}.
        NeurIPS 2020. \href{https://arxiv.org/abs/2005.11401}{arXiv:2005.11401}.

  \item Reimers, N., \& Gurevych, I. (2019).
        \textit{Sentence-BERT: Sentence Embeddings using Siamese BERT-Networks}.
        EMNLP 2019. \href{https://arxiv.org/abs/1908.10084}{arXiv:1908.10084}.

  \item OpenAI. (2024).
        \textit{GPT-4o Technical Report}.
        \href{https://openai.com/research/gpt-4o}{openai.com/research/gpt-4o}.

  \item Chromadb. (2024).
        \textit{ChromaDB: The AI-Native Open-Source Embedding Database}.
        \href{https://github.com/chroma-core/chroma}{github.com/chroma-core/chroma}.

  \item Yao, S., et al. (2023).
        \textit{ReAct: Synergizing Reasoning and Acting in Language Models}.
        ICLR 2023. \href{https://arxiv.org/abs/2210.03629}{arXiv:2210.03629}.

  \item Wang, L., et al. (2024).
        \textit{A Survey on Large Language Model based Autonomous Agents}.
        \textit{Frontiers of Computer Science}.
        \href{https://arxiv.org/abs/2308.11432}{arXiv:2308.11432}.

  \item Peng, B., et al. (2023).
        \textit{Check Your Facts and Try Again: Improving LLMs with External Knowledge and Automated Feedback}.
        \href{https://arxiv.org/abs/2302.12813}{arXiv:2302.12813}.

  \item Gao, Y., et al. (2024).
        \textit{Retrieval-Augmented Generation for Large Language Models: A Survey}.
        \href{https://arxiv.org/abs/2312.10997}{arXiv:2312.10997}.

  \item Hu, E., et al. (2022).
        \textit{LoRA: Low-Rank Adaptation of Large Language Models}.
        ICLR 2022. \href{https://arxiv.org/abs/2106.09685}{arXiv:2106.09685}.

  \item FastAPI. (2024).
        \textit{FastAPI Framework Documentation}.
        \href{https://fastapi.tiangolo.com}{fastapi.tiangolo.com}.
\end{enumerate}

\end{document}
